% WHICH CHARACTER NAMES AN ARGUMENT.
%
% A macro's parameter text may use more than one parameter character.
% \tracingmacros names each argument with the one THAT parameter was
% written with:
%
%   \def\a #1#2U3U4{...}   \a ABCD
%
%   \a #1#2U3U4->[U1|U2|U3|U4]
%   #1<-A
%   #2<-B
%   U3<-C
%   U4<-D
%
% The BODY, though, is written with the LAST of them throughout -- the
% reference carries one character through the whole list and the parameter
% text keeps setting it, so what is left when the body is reached is the
% last parameter's.
%
% This engine named every argument with the body's character. trip line 366
% defines \lo with the parameter text #1#2U3#4#5#6#7#8#989{ and line 367
% calls it.
\catcode`\{=1 \catcode`\}=2 \catcode`\#=6 \catcode`\U=6
\tracingonline=1 \tracingmacros=2
\message{[A a macro whose parameters use two different characters]}
\def\a #1#2U3U4{[#1|#2|U3|U4]}
\a ABCD
\message{[B and what \show makes of it]}
\show\a
\message{[C one that starts with the other character]}
\def\b U1#2{[U1|#2]}
\b EF
\show\b
\message{[done]}
\end
